% Generated by tools/docs/gen_fact_tables.py from rtl/soc/axi4_crossbar.sv and docs/manual/common/facts/bus.yaml. Do not edit.
{\small
\begin{longtable}{@{}>{\raggedright\arraybackslash}p{31mm}>{\raggedright\arraybackslash}p{9mm}>{\raggedright\arraybackslash}p{19mm}>{\raggedright\arraybackslash}p{22mm}>{\raggedright\arraybackslash}p{58mm}@{}}
\caption{Signals of \code{axi4_crossbar}}\label{tab:tm:mod:axi4-crossbar:signals}\\
\toprule
Signal & Dir & Width & Clock domain & Description\\
\midrule
\endfirsthead
\multicolumn{5}{@{}l}{\emph{Table \thetable\ continued}}\\
\toprule
Signal & Dir & Width & Clock domain & Description\\
\midrule
\endhead
\bottomrule
\endlastfoot
\code{clk} & in & 1 & clock & Clock of the crossbar. In the SoC it is \code{core_clk}.\\
\code{rst_n} & in & 1 & \code{clk} & Active-low synchronous reset. In the SoC it is \code{core_rst_n}.\\
\code{m_awaddr} & in & \code{[N_MASTERS-1:0][AW-1:0]} (96) & \code{clk} & Write address of the master of this index. The crossbar decodes it in the cycle that the engine accepts the address.\\
\code{m_awlen} & in & \code{[N_MASTERS-1:0][LENW-1:0]} (24) & \code{clk} & Burst length minus 1 of the write of the master of this index. The crossbar passes it to the slave.\\
\code{m_awsize} & in & \code{[N_MASTERS-1:0][2:0]} (9) & \code{clk} & Size code of a beat of the write of the master of this index. The crossbar passes it to the slave.\\
\code{m_awburst} & in & \code{[N_MASTERS-1:0][1:0]} (6) & \code{clk} & Burst type of the write of the master of this index. The crossbar passes it to the slave.\\
\code{m_awvalid} & in & \code{N_MASTERS} (3) & \code{clk} & The the master of this index offers a write address.\\
\code{m_awready} & out & \code{N_MASTERS} (3) & \code{clk} & The crossbar takes the write address of the master of this index. See the text about the second ready window.\\
\code{m_wdata} & in & \code{[N_MASTERS-1:0][DW-1:0]} (96) & \code{clk} & Write data of the master of this index.\\
\code{m_wstrb} & in & \code{[N_MASTERS-1:0][SW-1:0]} (12) & \code{clk} & Write strobes of the master of this index.\\
\code{m_wlast} & in & \code{N_MASTERS} (3) & \code{clk} & The beat is the last beat of the write of the master of this index. The crossbar ends the data phase on this bit.\\
\code{m_wvalid} & in & \code{N_MASTERS} (3) & \code{clk} & The the master of this index offers a write data beat.\\
\code{m_wready} & out & \code{N_MASTERS} (3) & \code{clk} & The crossbar takes the write data beat of the master of this index. It is high only while the write engine of the granted slave is in its data phase, or while the decode-error engine of the master sinks data.\\
\code{m_bresp} & out & \code{[N_MASTERS-1:0][1:0]} (6) & \code{clk} & Write response to the master of this index. The response of the slave, or DECERR from the decode-error engine.\\
\code{m_bvalid} & out & \code{N_MASTERS} (3) & \code{clk} & The crossbar offers the write response to the master of this index.\\
\code{m_bready} & in & \code{N_MASTERS} (3) & \code{clk} & The the master of this index takes the write response.\\
\code{m_araddr} & in & \code{[N_MASTERS-1:0][AW-1:0]} (96) & \code{clk} & Read address of the master of this index. The crossbar decodes it in the cycle that the engine accepts the address.\\
\code{m_arlen} & in & \code{[N_MASTERS-1:0][LENW-1:0]} (24) & \code{clk} & Burst length minus 1 of the read of the master of this index. The crossbar passes it to the slave.\\
\code{m_arsize} & in & \code{[N_MASTERS-1:0][2:0]} (9) & \code{clk} & Size code of a beat of the read of the master of this index. The crossbar passes it to the slave.\\
\code{m_arburst} & in & \code{[N_MASTERS-1:0][1:0]} (6) & \code{clk} & Burst type of the read of the master of this index. The crossbar passes it to the slave.\\
\code{m_arvalid} & in & \code{N_MASTERS} (3) & \code{clk} & The the master of this index offers a read address.\\
\code{m_arready} & out & \code{N_MASTERS} (3) & \code{clk} & The crossbar takes the read address of the master of this index. See the text about the second ready window.\\
\code{m_rdata} & out & \code{[N_MASTERS-1:0][DW-1:0]} (96) & \code{clk} & Read data to the master of this index. It is 0 for a decode-error beat.\\
\code{m_rresp} & out & \code{[N_MASTERS-1:0][1:0]} (6) & \code{clk} & Read response to the master of this index. The response of the slave, or DECERR from the decode-error engine.\\
\code{m_rlast} & out & \code{N_MASTERS} (3) & \code{clk} & The beat is the last beat of the read to the master of this index.\\
\code{m_rvalid} & out & \code{N_MASTERS} (3) & \code{clk} & The crossbar offers a read data beat to the master of this index.\\
\code{m_rready} & in & \code{N_MASTERS} (3) & \code{clk} & The the master of this index takes the read data beat.\\
\code{s_awid} & out & \code{[N_SLAVES-1:0][IW-1:0]} (8) & \code{clk} & ID of the write: the index of the granted master.\\
\code{s_awaddr} & out & \code{[N_SLAVES-1:0][AW-1:0]} (64) & \code{clk} & Address of the write, taken from the register that captured it.\\
\code{s_awlen} & out & \code{[N_SLAVES-1:0][LENW-1:0]} (16) & \code{clk} & Burst length minus 1 of the write.\\
\code{s_awsize} & out & \code{[N_SLAVES-1:0][2:0]} (6) & \code{clk} & Size code of a beat of the write.\\
\code{s_awburst} & out & \code{[N_SLAVES-1:0][1:0]} (4) & \code{clk} & Burst type of the write.\\
\code{s_awvalid} & out & \code{N_SLAVES} (2) & \code{clk} & The crossbar offers the write address to the slave of this index. It is high in the address state of the write engine.\\
\code{s_awready} & in & \code{N_SLAVES} (2) & \code{clk} & Slave \code{s} takes the write address.\\
\code{s_wdata} & out & \code{[N_SLAVES-1:0][DW-1:0]} (64) & \code{clk} & Write data of the granted master. It is 0 outside the data state.\\
\code{s_wstrb} & out & \code{[N_SLAVES-1:0][SW-1:0]} (8) & \code{clk} & Write strobes of the granted master. They are 0 outside the data state.\\
\code{s_wlast} & out & \code{N_SLAVES} (2) & \code{clk} & Last-beat flag of the granted master. It is 0 outside the data state.\\
\code{s_wvalid} & out & \code{N_SLAVES} (2) & \code{clk} & The write data of the granted master is valid. It is 0 outside the data state.\\
\code{s_wready} & in & \code{N_SLAVES} (2) & \code{clk} & Slave \code{s} takes the write data beat.\\
\code{s_bid} & in & \code{[N_SLAVES-1:0][IW-1:0]} (8) & \code{clk} & Write response ID from the slave of this index. The crossbar does not use it.\\
\code{s_bresp} & in & \code{[N_SLAVES-1:0][1:0]} (4) & \code{clk} & Write response of the slave of this index.\\
\code{s_bvalid} & in & \code{N_SLAVES} (2) & \code{clk} & Slave \code{s} offers the write response.\\
\code{s_bready} & out & \code{N_SLAVES} (2) & \code{clk} & The crossbar takes the write response. It follows \code{m_bready} of the granted master in the response state.\\
\code{s_arid} & out & \code{[N_SLAVES-1:0][IW-1:0]} (8) & \code{clk} & ID of the read: the index of the granted master.\\
\code{s_araddr} & out & \code{[N_SLAVES-1:0][AW-1:0]} (64) & \code{clk} & Address of the read, taken from the register that captured it.\\
\code{s_arlen} & out & \code{[N_SLAVES-1:0][LENW-1:0]} (16) & \code{clk} & Burst length minus 1 of the read.\\
\code{s_arsize} & out & \code{[N_SLAVES-1:0][2:0]} (6) & \code{clk} & Size code of a beat of the read.\\
\code{s_arburst} & out & \code{[N_SLAVES-1:0][1:0]} (4) & \code{clk} & Burst type of the read.\\
\code{s_arvalid} & out & \code{N_SLAVES} (2) & \code{clk} & The crossbar offers the read address to the slave of this index. It is high in the address state of the read engine.\\
\code{s_arready} & in & \code{N_SLAVES} (2) & \code{clk} & Slave \code{s} takes the read address.\\
\code{s_rid} & in & \code{[N_SLAVES-1:0][IW-1:0]} (8) & \code{clk} & Read data ID from the slave of this index. The crossbar does not use it.\\
\code{s_rdata} & in & \code{[N_SLAVES-1:0][DW-1:0]} (64) & \code{clk} & Read data of the slave of this index.\\
\code{s_rresp} & in & \code{[N_SLAVES-1:0][1:0]} (4) & \code{clk} & Read response of the slave of this index.\\
\code{s_rlast} & in & \code{N_SLAVES} (2) & \code{clk} & Slave \code{s} marks the last read beat. The read engine ends on this bit.\\
\code{s_rvalid} & in & \code{N_SLAVES} (2) & \code{clk} & Slave \code{s} offers a read data beat.\\
\code{s_rready} & out & \code{N_SLAVES} (2) & \code{clk} & The crossbar takes the read data beat. It follows \code{m_rready} of the granted master in the data state.\\
\end{longtable}
}
